\label{sec:problem_setting}

The application target is local railway-domain bilingual assistance rather than
general-purpose online chatbot service. The intended users are railway
engineering students, laboratory researchers, and technical staff who need
fast access to bilingual terminology and operation-and-maintenance knowledge.
In this setting, the model receives a short technical prompt and is expected to
return a concise answer without long reasoning traces, unrelated explanation,
or unsafe operational advice.

The input space contains two common task families. The first is terminology
translation, where the model maps a railway technical expression between
Chinese and English. The second is regulation-oriented translation or
question-answering, where the prompt contains a sentence or paragraph from
traction power supply operation-and-maintenance material. The output is a
single target-language answer or domain-specific phrase. These tasks are small
in context length, but difficult for generic models because many answers depend
on specialized railway vocabulary and conventional wording.

The deployment constraint is fixed to a single NVIDIA RTX 3090 GPU with 24GB of
VRAM. This constraint reflects a realistic laboratory workstation and rules out
multi-GPU serving, full-parameter fine-tuning of larger models, and expensive
hyperparameter searches. Therefore, the model selection problem is
multi-objective: the selected model must be accurate enough for technical use,
fit reliably into memory, answer at usable latency, and support a low-cost
adaptation path.

The current dataset contains 22,164 training samples, 2,770 validation samples,
and 2,770 test samples. A 1\% smoke subset is used to validate the pipeline
before longer runs. The present result export includes public benchmark scores
for six completed model-precision combinations and a one-sample Domain QA
sanity check. This means the draft reports mature public benchmark and
efficiency measurements, while the domain QA and QLoRA results should be read
as pipeline-completeness evidence until the full 2,770-sample domain test
evaluation is rerun and exported.
